%% LyX 2.4.4 created this file.  For more info, see https://www.lyx.org/.
%% Do not edit unless you really know what you are doing.
\documentclass[english]{article}
\usepackage[T1]{fontenc}
\usepackage[utf8]{inputenc}
\setcounter{secnumdepth}{-2}
\setlength{\parindent}{0bp}
\usepackage{amsmath}
\usepackage{amsthm}
\usepackage{amssymb}

\makeatletter
%%%%%%%%%%%%%%%%%%%%%%%%%%%%%% Textclass specific LaTeX commands.
\theoremstyle{definition}
\newtheorem*{problem*}{\protect\problemname}

\@ifundefined{date}{}{\date{}}
%%%%%%%%%%%%%%%%%%%%%%%%%%%%%% User specified LaTeX commands.
\usepackage{graphicx}
\usepackage{geometry}
\newcommand{\limi}{\underset{n\rightarrow\infty}{\lim}}
\newcommand{\limxz}{\underset{x\rightarrow x_{0}}{\lim}}
\newcommand{\limfxz}{\underset{x\rightarrow x_{0}}{\lim}f(x)}
\newcommand{\limfxm}{\underset{x\rightarrow x_{0}^{-}}{\lim}f(x)}
\newcommand{\limfxp}{\underset{x\rightarrow x_{0}^{+}}{\lim}f(x)}
\newcommand{\limxm}{\underset{x\rightarrow x_{0}^{-}}{\lim}}
\newcommand{\limxp}{\underset{x\rightarrow x_{0}^{+}}{\lim}}
\newcommand{\sqrm}{M_{m\times m}(\mathbb{F})}
\newcommand{\seqa}{(a_{n})_{n=1}^{\infty}}
\newcommand{\seqx}{(x_{n})_{n=1}^{\infty}}
\newcommand{\funcdr}{f:D\rightarrow\mathbb{R}}
\newcommand{\der}{\limxz\frac{f(x)-f(x_{0})}{x-x_{0}}}
\usepackage[all]{pdfprivacy}

\makeatother

\usepackage{babel}
\providecommand{\problemname}{Problem}

\begin{document}
\title{{\Large Semester 2026A, Exam B, Q3\vspace{-2em}}}
\maketitle
\begin{problem*}
Let $V$ be a vector space over $\mathbb{R}$ and let $(v_{1},v_{2},v_{3})$
be a linearly independent sequence of vectors from $V$. Determine
for which values $a\in\mathbb{R}$ the sequence 
\[
S=(v_{1}+av_{2},\;v_{2}+av_{3},\;v_{3}+av_{1})
\]
is also linearly independent. Prove your answer.
\end{problem*}
\medskip{}

\rule[0.5ex]{1\columnwidth}{1pt}

\medskip{}

\begin{proof}
Let us recall the theorem which states that a sequence of vectors
is linearly independent if and only if the only way to reach $0$
as a linear combination of it is by taking all of the scalars as zero. 

Let us examine the following equation 

\[
\alpha_{1}\left(v_{1}+av_{2}\right)+\alpha_{2}\left(v_{2}+av_{3}\right)+\alpha_{3}\left(v_{3}+av_{1}\right)=0
\]
\[
\alpha_{1}v_{1}+\alpha_{1}av_{2}+\alpha_{2}v_{2}+\alpha_{2}av_{3}+\alpha_{3}v_{3}+\alpha_{3}av_{1}=0
\]
\[
\left(\alpha_{1}+\alpha_{3}a\right)v_{1}+\left(\alpha_{1}a+\alpha_{2}\right)v_{2}+\left(\alpha_{2}a+\alpha_{3}\right)v_{3}=0
\]

And from the same theorem, since $(v_{1},v_{2},v_{3})$ is linearly
independent, we know that 
\[
\begin{cases}
\left(\alpha_{1}+\alpha_{3}a\right)=0\\
\left(\alpha_{1}a+\alpha_{2}\right)=0\\
\left(\alpha_{2}a+\alpha_{3}\right)=0
\end{cases}
\]

We will use this linear system to determine which $a$ values give
us $\alpha_{1}=\alpha_{2}=\alpha_{3}=0$, i.e. that the only solution
is the trivial solution.

\[
\begin{pmatrix}\begin{array}{ccc|c}
1 & 0 & a & 0\\
a & 1 & 0 & 0\\
0 & a & 1 & 0
\end{array}\end{pmatrix}\overset{R_{2}\rightarrow R_{2}-aR_{1}}{\Longrightarrow}\begin{pmatrix}\begin{array}{ccc|c}
1 & 0 & a & 0\\
0 & 1 & -a^{2} & 0\\
0 & a & 1 & 0
\end{array}\end{pmatrix}\overset{R_{3}\rightarrow R_{3}-aR_{2}}{\Longrightarrow}\begin{pmatrix}\begin{array}{ccc|c}
1 & 0 & a & 0\\
0 & 1 & -a^{2} & 0\\
0 & 0 & 1+a^{3} & 0
\end{array}\end{pmatrix}
\]

$\left(1+a^{3}\right)\alpha_{3}=0$ means $\alpha_{3}=0$ or $a^{3}=-1\iff a=-1$.

Considering this, $\alpha_{2}-a^{2}\alpha_{3}=0$ means that if $\alpha_{3}=0$
then $\alpha_{2}=0$ and if $a=-1$ then $\alpha_{2}=\alpha_{3}$.

Similarly, $\alpha_{1}+a\alpha_{3}=0$ means that if $\alpha_{3}=0$
then $\alpha_{1}=0$ and if $a=-1$ then $\alpha_{1}=\alpha_{3}$. 

Which means that if $a\neq-1$ the set of solutions is $\left\{ \begin{pmatrix}0\\
0\\
0
\end{pmatrix}\right\} $ and if $a=-1$ it's $\left\{ \begin{pmatrix}x\\
x\\
x
\end{pmatrix}\mid x\in\mathbb{R}\right\} $, which means that the sequence is linearly independent for $a\in\mathbb{R}\setminus\{-1\}$.
\end{proof}
\begin{quote}
This is an exam problem I translated and solved. The original exam
was written by the Linear Algebra 1 course lecturers at HUJI.
\end{quote}

\end{document}
